% from ega1-5-fr.tex
\oldpage[I]{129}
phisme $\theta_x^\sharp:\mathscr{O}_{\psi(x)}\to\mathscr{O}_x$ applique tout
élément nilpotent de $\mathscr{O}_{\psi(x)}$ sur un élément nilpotent de
$\mathscr{O}_x$~; par passage aux quotients, on déduit donc de $\theta^\sharp$
un homomorphisme
\[
  \omega:\psi^*(\mathscr{O}_Y/\mathscr{N}_Y)
  \longrightarrow \mathscr{O}_X/\mathscr{N}_X
\]
il est clair que pour tout $x\in X$,
$\omega_x:\mathscr{O}_{\psi(x)}/\mathscr{N}_{\psi(x)}
\to\mathscr{O}_x/\mathscr{N}_x$ est un homomorphisme local, donc
$(\psi,\omega^b)$ est un morphisme de préschémas
$X_{\mathrm{red}}\to Y_{\mathrm{red}}$, que nous noterons $f_{\mathrm{red}}$
et appellerons le morphisme \emph{réduit} associé à $f$. Il est immédiat que
pour deux morphismes $f:X\to Y$, $g:Y\to Z$, on a
$(g\circ f)_{\mathrm{red}}=g_{\mathrm{red}}\circ f_{\mathrm{red}}$, donc on a
défini $X_{\mathrm{red}}$ comme \emph{foncteur covariant} en $X$.

La définition précédente montre que le diagramme
\[
  \xymatrix{
    X_{\mathrm{red}}\ar[r]^{f_{\mathrm{red}}}\ar[d] &
    Y_{\mathrm{red}}\ar[d]\\
    X\ar[r]_f & Y
  }
\]
est commutatif, les flèches verticales étant les morphismes d'injection~; en
d'autres termes, $X_{\mathrm{red}}\to X$ est un morphisme \emph{fonctoriel}.
On notera en particulier que si $X$ est \emph{réduit}, tout morphisme
$f:X\to Y$ se factorise en
$X\xrightarrow{f_{\mathrm{red}}}Y_{\mathrm{red}}\to Y$~; en d'autres termes,
$f$ est \emph{majoré} par le morphisme d'injection
$Y_{\mathrm{red}}\to Y$.
\end{env}

\begin{proposition}[5.1.6]
\label{I.5.1.6-fr}
Soit $f:X\to Y$ un morphisme~; si $f$ est surjectif (resp. radiciel, une
immersion, une immersion fermée, une immersion ouverte, une immersion locale,
un isomorphisme local), il en est de même de $f_{\mathrm{red}}$. Inversement,
si $f_{\mathrm{red}}$ est surjectif (resp. radiciel), il en est de même de
$f$.
\end{proposition}

La proposition est triviale si $f$ est surjectif~; si $f$ est radiciel, elle
résulte de ce que pour tout $x\in X$, le corps $k(x)$ est le même pour les
préschémas $X$ et $X_{\mathrm{red}}$
(\hyperref[I.3.5.8-fr]{3.5.8}). Enfin, si $f=(\psi,\theta)$ est une immersion,
une immersion fermée ou une immersion locale (resp. une immersion ouverte, ou
un isomorphisme local), la proposition résulte de ce que si
$\theta_x^\sharp$ est surjectif (resp. bijectif), il en est de même de
l'homomorphisme obtenu en passant aux quotients par les nilradicaux de
$\mathscr{O}_{\psi(x)}$ et $\mathscr{O}_x$
(\hyperref[I.5.1.2-fr]{5.1.2} et \hyperref[I.4.2.2-fr]{4.2.2})
(cf. (\hyperref[I.5.5.12-fr]{5.5.12})).

\begin{proposition}[5.1.7]
\label{I.5.1.7-fr}
Si $X$, $Y$ sont deux $S$-préschémas, les préschémas
$X_{\mathrm{red}}\times_{S_{\mathrm{red}}}Y_{\mathrm{red}}$ et
$X_{\mathrm{red}}\times_S Y_{\mathrm{red}}$ sont identiques, et
s'identifient canoniquement à un sous-préschéma de $X\times_S Y$ ayant même
espace sous-jacent que ce produit.
\end{proposition}

L'identification canonique de $X_{\mathrm{red}}\times_S Y_{\mathrm{red}}$ à
un sous-préschéma de $X\times_S Y$ ayant même espace sous-jacent résulte de
(\hyperref[I.4.3.1-fr]{4.3.1}). D'autre part, si $\varphi$ et $\psi$ sont les
morphismes structuraux $X_{\mathrm{red}}\to S$, $Y_{\mathrm{red}}\to S$, ils
se factorisent par $S_{\mathrm{red}}$
(\hyperref[I.5.1.5-fr]{5.1.5}), et comme $S_{\mathrm{red}}\to S$ est un
monomorphisme, la première assertion résulte de
(\hyperref[I.3.2.4-fr]{3.2.4}).

\begin{corollary}[5.1.8]
\label{I.5.1.8-fr}
Les préschémas $(X\times_S Y)_{\mathrm{red}}$ et
$(X_{\mathrm{red}}\times_{S_{\mathrm{red}}}Y_{\mathrm{red}})_{\mathrm{red}}$
s'identifient canoniquement.
\end{corollary}

Cela résulte de (\hyperref[I.5.1.2-fr]{5.1.2} et
\hyperref[I.5.1.7-fr]{5.1.7}).

On notera que si $X$ et $Y$ sont des préschémas réduits, il n'en est pas
nécessairement de même de $X\times_S Y$, car le produit tensoriel de deux
algèbres réduites peut avoir des éléments nilpotents.

\begin{proposition}[5.1.9]
\label{I.5.1.9-fr}
